% Generated draft: requires mathematical and source review.
\documentclass[11pt]{article}
\usepackage[margin=25mm]{geometry}
\usepackage{fontspec}
\setmainfont{Latin Modern Roman}
\usepackage{amsmath,amssymb,mathtools,unicode-math}
\setmathfont{Latin Modern Math}
\usepackage{enumitem,needspace,xurl,hyperref}
\setlength{\parindent}{0pt}
\setlength{\parskip}{5pt}
\setlength{\emergencystretch}{3em}
\newcommand{\N}{\mathbb{N}}
\newcommand{\Z}{\mathbb{Z}}
\newcommand{\Q}{\mathbb{Q}}
\newcommand{\R}{\mathbb{R}}
\newcommand{\C}{\mathbb{C}}
\newcommand{\F}{\mathbb{F}}
\newcommand{\A}{\mathbb{A}}
\newcommand{\PP}{\mathbb{P}}
\newcommand{\E}{\mathbb{E}}
\newcommand{\eps}{\varepsilon}
\newcommand{\dd}{\,\mathrm{d}}
\begin{document}
\section*{UnsolvedMath ID: 30000042}
\textbf{Source problem code:} catalogue entry
\subsection{Definitions and mathematical statement}
Let \(H\) be a real or complex Hilbert space and \(M\subset H\) a closed subspace. For an index set \(I\), a frame \(F=(f_i)_{i\in I}\) for \(M\) satisfies \(f_i\in M\) and
\[
A\|x\|^2\leq\sum_{i\in I}|\langle x,f_i\rangle|^2\leq B\|x\|^2
\quad(x\in M),\qquad 0<A\leq B<\infty.
\]
Its frame operator on \(M\) is \(S_Fx=\sum_i\langle x,f_i\rangle f_i\), with the inner-product convention linear in the first argument. A normalized tight frame means a Parseval frame, for which the displayed sum equals \(\|x\|^2\). The canonical associated normalized tight frame is
\[
F^{\mathrm{can}}=(S_F^{-1/2}f_i)_{i\in I}.
\]
This is a frame for \(M\); when \(M\ne H\), the inverse is taken on \(M\), not on all of \(H\).

The source asks for metrics \(d\) on a specified space of frames, and appropriate admissible normalized tight frames \(\mathcal P\), such that
\[
d(F,F^{\mathrm{can}})=\inf_{G\in\mathcal P}d(F,G)
\]
for every input frame in the intended class. The equality permits several closest frames. It also asks for the geometry of the subspaces supporting closest frames. If competitors may span different closed subspaces of \(H\), a precise notation for that object is
\[
\mathcal C_d(F)=\{N\subset H:N\text{ closed},\ \exists G\text{ a Parseval frame for }N,
\ G\in\mathcal P,\ d(F,G)=\inf_{K\in\mathcal P}d(F,K)\}.
\]
Competitors must have compatible indexing whenever \(d\) compares individual frame vectors. If all competitors are required to span the fixed \(M\), the minimization problem changes and the subspace set above has no variation. The source does not choose between fixed-subspace and varying-subspace competitors, identify the admissible metrics, or specify finite versus infinite dimensions and indices. Those data are needed for a fully quantified classification theorem. In infinite settings, a distance given by a sum of squared vector differences also needs finiteness assumptions and may not define a finite metric on all frames.

The supplied assessment dated 2026-08-21 reports nearest-tight-frame results for important group-generated classes, but says that the question for arbitrary frames, subspaces, and distances is not settled in full generality. No particular metric or theorem for those classes is reproduced in the record.
\subsection{Short English statement}
Identify distances under which the canonical Parseval frame is a nearest Parseval approximation, and describe the subspaces on which all nearest competitors lie. The admissible distances, frame classes, and whether competitors may change subspace remain unspecified.
\subsection{Sources}
\begin{itemize}
\item Local UnsolvedMath registry, permanent ID 30000042, catalogue entry, \emph{Nearest Normalized Tight Frames}.
\item \emph{Mini-Workshop: Wavelets and Frames} (2004), Oberwolfach Reports, pp. 479--544, DOI \url{https://doi.org/10.4171/owr/2004/10}.
\item Supplied reference labeled ``Han, Gabor and wavelet frames'': \url{https://doi.org/10.1090/S0002-9947-03-03047-2}.
\item Supplied reference labeled ``Balan, symmetric approximation'': \url{https://arxiv.org/abs/math/9812052}.
\item Supplied literature triage checked 2026-08-21, reporting partial results for group-generated frames; full theorem statements are not supplied.
\end{itemize}
\end{document}
